% GENERATED FILE. Edit canonical lessons, then run npm run generate:books.
% canonical-source-hash: fnv1a64:f43ec02b79b97709
% canonical-lessons: JA-C255-itta-hou-ga-ii, JA-C255-dakara, JA-C255-desu-kara, JA-R255-advice-and-so

\chapter{Advice, and So}
\label{ch:ja-advice-and-so}
\hlchaptermodality{255}

\begin{chapteropening}
\textbf{By the end of this chapter:} \emph{I can give advice for and against (yasunda hō ga ii desu; nomanai hō ga ii desu), and say what follows from something with dakara, desu kara and sorede.}

Complete JA-R255-advice-and-so: give advice and say what follows from it.
\end{chapteropening}

\section[\texorpdfstring{\ja{いった} \ja{ほうが} … (itta hō ga ii desu)}{いった ほうが … (itta hō ga ii desu)}]{\ja{いった} \ja{ほうが} \ja{いいです} (itta hō ga ii desu) — you had better}

\label{lesson:JA-C255-itta-hou-ga-ii}



\begin{quote}
Say, plainly, \textquotedblleft{}I rested\textquotedblright{} and \textquotedblleft{}I won't drink\textquotedblright{}. (\emph{yasunda}, \emph{nomanai}.) Then say \textquotedblleft{}medicine\textquotedblright{}. (\emph{kusuri}.)
\end{quote}

\subsection*{You'll want to know: -\ja{た} \ja{ほうが} \ja{いいです}}
\begin{quote}
\textbf{\ja{やすんだ} \ja{ほうが} \ja{いいです。}}
\end{quote}

\begin{quote}
\emph{yasunda hō ga ii desu.} — \textquotedblleft{}You had better rest.\textquotedblright{}
\end{quote}

\textbf{\ja{ほう}} (\emph{hō}) is \textquotedblleft{}side, way\textquotedblright{}: \textquotedblleft{}the side of having rested is good\textquotedblright{}. The plain past goes before it even when the advice is about now or later.

\begin{quote}
\textbf{\ja{くすりを} \ja{のんだ} \ja{ほうが} \ja{いいです。}}
\end{quote}

\begin{quote}
\emph{kusuri o nonda hō ga ii desu.} — \textquotedblleft{}You had better take some medicine.\textquotedblright{}
\end{quote}

\begin{quote}
\textbf{\ja{びょういんに} \ja{いった} \ja{ほうが} \ja{いいです。}}
\end{quote}

\begin{quote}
\emph{byōin ni itta hō ga ii desu.} — \textquotedblleft{}You had better go to the hospital.\textquotedblright{}
\end{quote}

\subsection*{You'll want to know: -\ja{ない} \ja{ほうが} \ja{いいです}}
Advice against takes the plain \textquotedblleft{}not\textquotedblright{}:

\begin{quote}
\textbf{\ja{おさけを} \ja{のまない} \ja{ほうが} \ja{いいです。}}
\end{quote}

\begin{quote}
\emph{osake o nomanai hō ga ii desu.} — \textquotedblleft{}You had better not drink alcohol.\textquotedblright{}
\end{quote}

\begin{quote}
\textbf{\ja{あぶないから、いかない} \ja{ほうが} \ja{いいです。}}
\end{quote}

\begin{quote}
\emph{abunai kara, ikanai hō ga ii desu.} — \textquotedblleft{}It is dangerous, so you had better not go.\textquotedblright{}
\end{quote}

This is advice. Saying that something must be done needs another form, which comes later.

\subsection*{Your turn}
Say these aloud:

\begin{itemize}
\raggedright
  \item \textquotedblleft{}You had better not drink alcohol\textquotedblright{} — \emph{osake o nomanai hō ga ii desu}
  \item \textquotedblleft{}You had better not drink coffee\textquotedblright{} — \emph{kōhī o nomanai hō ga ii desu}
  \item \textquotedblleft{}You had better not drink cold water\textquotedblright{} — \emph{tsumetai mizu o nomanai hō ga ii desu}
\end{itemize}

\begin{tcolorbox}[breakable,colback=teal!4,colframe=teal!35!black,title={Before you move on}]
\textquotedblleft{}You had better rest\textquotedblright{}? (\textbf{\ja{やすんだ} \ja{ほうが} \ja{いいです}}.) \textquotedblleft{}You had better not go\textquotedblright{}? (\textbf{\ja{いかない} \ja{ほうが} \ja{いいです}}.) Which form goes before \textbf{\ja{ほうが}} for advice for? (The plain past.)
\end{tcolorbox}
\section[\texorpdfstring{\ja{だから} (dakara)}{だから (dakara)}]{\ja{だから} (dakara) — so; that is why}

\label{lesson:JA-C255-dakara}



\begin{quote}
Before the new word: say \textquotedblleft{}it is cold, so I won't go\textquotedblright{}. (\emph{samui kara, ikimasen}.) Then \textquotedblleft{}you had better not go\textquotedblright{}. (\emph{ikanai hō ga ii desu}.)
\end{quote}

\subsection*{You'll want to know: \ja{だから}}
\textbf{\ja{だから}} — \emph{dakara} — \textquotedblleft{}so\textquotedblright{}, \textquotedblleft{}that is why\textquotedblright{}. It opens a new sentence with the result, after a sentence that gave the reason:

\begin{quote}
\textbf{\ja{あめが} \ja{ふっています。}}
\end{quote}

\begin{quote}
\emph{ame ga futte imasu.} — \textquotedblleft{}It is raining.\textquotedblright{}
\end{quote}

\begin{quote}
\textbf{\ja{だから、いきません。}}
\end{quote}

\begin{quote}
\emph{dakara, ikimasen.} — \textquotedblleft{}So I won't go.\textquotedblright{}
\end{quote}

It is plain, the word for a friend; the next lesson gives its polite partner.

\begin{cousinweb}
\textbf{\ja{だから}} (\emph{dakara}) is a native Japanese word (wago), inherited rather than borrowed from Chinese: \textbf{\ja{だ}} (\emph{da}), the plain form of \textbf{\ja{です}}, and \textbf{\ja{から}} (\emph{kara}), \textquotedblleft{}because\textquotedblright{}. JMdict lists it as a conjunction, \textquotedblleft{}so; therefore\textquotedblright{}, written in kana (JMdict).
\end{cousinweb}

\subsection*{Your turn}
Say these aloud:

\begin{itemize}
\raggedright
  \item \emph{dakara}
  \item \textquotedblleft{}So I won't go\textquotedblright{} — \emph{dakara, ikimasen}
  \item \textquotedblleft{}So I will rest\textquotedblright{} — \emph{dakara, yasumimasu}
\end{itemize}

\begin{tcolorbox}[breakable,colback=teal!4,colframe=teal!35!black,title={Before you move on}]
\textquotedblleft{}So\textquotedblright{}, opening a sentence, to a friend? (\textbf{\ja{だから}}, \emph{dakara}.) Which comes first, the reason or \textbf{\ja{だから}}? (The reason, in the sentence before.)
\end{tcolorbox}
\section[\texorpdfstring{\ja{ですから} / \ja{それで} (desu kara / sore de)}{ですから / それで (desu kara / sore de)}]{\ja{ですから} / \ja{それで} (desu kara / sore de) — so, and so}

\label{lesson:JA-C255-desu-kara}



\begin{quote}
Say \textquotedblleft{}so I won't go\textquotedblright{}, to a friend. (\emph{dakara, ikimasen}.) Then say \textquotedblleft{}a holiday\textquotedblright{}. (\emph{yasumi}.)
\end{quote}

\subsection*{You'll want to know: \ja{ですから}}
To a teacher, in a shop or at work, \textbf{\ja{だから}} becomes \textbf{\ja{ですから}} (\emph{desu kara}), with \textbf{\ja{です}} in place of \textbf{\ja{だ}}:

\begin{quote}
\textbf{\ja{あしたは} \ja{やすみです。}}
\end{quote}

\begin{quote}
\emph{ashita wa yasumi desu.} — \textquotedblleft{}Tomorrow is a holiday.\textquotedblright{}
\end{quote}

\begin{quote}
\textbf{\ja{ですから、みせは} \ja{しまっています。}}
\end{quote}

\begin{quote}
\emph{desu kara, mise wa shimatte imasu.} — \textquotedblleft{}So the shop is closed.\textquotedblright{}
\end{quote}

\subsection*{You'll want to know: \ja{それで}}
\textbf{\ja{それで}} (\emph{sore de}) is \textquotedblleft{}and so\textquotedblright{}: it tells what came of something that happened. It looks back at an event, where \textbf{\ja{だから}} and \textbf{\ja{ですから}} also give advice and decisions.

\begin{quote}
\textbf{\ja{でんしゃが} \ja{とまりました。}}
\end{quote}

\begin{quote}
\emph{densha ga tomarimashita.} — \textquotedblleft{}The train stopped.\textquotedblright{}
\end{quote}

\begin{quote}
\textbf{\ja{それで、おくれました。}}
\end{quote}

\begin{quote}
\emph{sore de, okuremashita.} — \textquotedblleft{}And so I was late.\textquotedblright{}
\end{quote}

\subsection*{Your turn}
Say these aloud:

\begin{itemize}
\raggedright
  \item \textquotedblleft{}And so I was late\textquotedblright{} — \emph{sore de, okuremashita}
  \item \textquotedblleft{}And so I went home\textquotedblright{} — \emph{sore de, kaerimashita}
  \item \textquotedblleft{}And so I walked\textquotedblright{} — \emph{sore de, arukimashita}
\end{itemize}

\begin{tcolorbox}[breakable,colback=teal!4,colframe=teal!35!black,title={Before you move on}]
\textquotedblleft{}So\textquotedblright{}, politely? (\textbf{\ja{ですから}}.) \textquotedblleft{}So\textquotedblright{}, to a friend? (\textbf{\ja{だから}}.) \textquotedblleft{}And so\textquotedblright{}, for what came of an event? (\textbf{\ja{それで}}.)
\end{tcolorbox}
\section[(review)]{Review — advice, and so}

\label{lesson:JA-R255-advice-and-so}



\begin{quote}
Say \textquotedblleft{}you had better rest\textquotedblright{}. (\emph{yasunda hō ga ii desu}.) Then \textquotedblleft{}so I won't go\textquotedblright{}, to a friend. (\emph{dakara, ikimasen}.)
\end{quote}

\subsection*{Your turn}
Say these aloud:

\begin{itemize}
\raggedright
  \item You had better take some medicine.
  \item You had better not drink cold water.
  \item It is raining. So I won't go. — to a friend
  \item The train stopped. And so I was late.
  \item I think you had better go to the hospital.
  \item I asked my teacher's opinion.
\end{itemize}

\begin{tcolorbox}[breakable,colback=teal!4,colframe=teal!35!black,title={Before you move on}]
Advice for? (\textbf{-\ja{た} \ja{ほうが} \ja{いいです}}.) Advice against? (\textbf{-\ja{ない} \ja{ほうが} \ja{いいです}}.) \textquotedblleft{}So\textquotedblright{}, to a friend and politely? (\textbf{\ja{だから}}, \textbf{\ja{ですから}}.) \textquotedblleft{}And so\textquotedblright{}, for what came of it? (\textbf{\ja{それで}}.)
\end{tcolorbox}
